\section{Corporate Database}
\label{sec:ch06-corporate-database}
\begin{definitionbox}{}
প্রতিষ্ঠানের বিভিন্ন বিভাগের কার্যক্রমে প্রয়োজনীয় সম্পর্কিত উপাত্ত সমন্বিতভাবে সংরক্ষণ ও ব্যবহারের ডেটাবেজকে কর্পোরেট ডেটাবেজ বলে।
\end{definitionbox}


Corporate Database ও Government Database পৃথক Data Model নয়; এগুলো ব্যবহারের ক্ষেত্রভিত্তিক নাম।


\subsubsection{গ্রাহকের একটি প্রশ্ন, তিন বিভাগের আলাদা তথ্য}
একজন গ্রাহক পণ্যের দাম দিয়েছেন, কিন্তু পণ্যটি এখনো পাননি। বিক্রয় বিভাগে (Sales) অর্ডারের তথ্য, হিসাব বিভাগে (Accounts) পরিশোধের তথ্য এবং পণ্য পাঠানোর বিভাগে (Delivery) সরবরাহের তথ্য আছে। এসব তথ্য মিলিয়ে দেখা না গেলে গ্রাহককে প্রতিটি বিভাগে একই অভিযোগ জানাতে হয়।

তাই অনুমতি অনুযায়ী সংশ্লিষ্ট তথ্য একসঙ্গে পাওয়ার ব্যবস্থা প্রয়োজন। CustomerID ও OrderID দিয়ে তথ্যগুলো মিলিয়ে সেবাকর্মী বুঝতে পারেন, কোন ধাপে কাজ থেমেছে। অভিযোগ সমাধানের অগ্রগতিও সংরক্ষণ করা যায়। দোকানের গ্রাহকসেবা, হাসপাতালের বিল ও সেবা এবং প্রতিষ্ঠানের উপস্থিতি ও বেতনের হিসাব সমন্বয়ে এ ধরনের ব্যবস্থা কাজে লাগে। এতে তথ্য সংরক্ষণের পাশাপাশি বিভিন্ন বিভাগের কাজ কীভাবে এগোচ্ছে, তাও বোঝা যায়।


বৃহৎ ব্যবসা, বিশ্ববিদ্যালয়, হাসপাতাল বা গবেষণা প্রতিষ্ঠানকে তাদের সংগঠিত কার্যক্রমের দিক থেকে Enterprise বলা যায়। এসব প্রতিষ্ঠানে একই তথ্য বিক্রয়, হিসাব, মজুদ ও ব্যবস্থাপনার কাজে লাগতে পারে।

\subsection{বিভাগভিত্তিক ব্যবহারের ক্ষেত্র}
\begin{enumerate}
\item \textbf{মানবসম্পদ ব্যবস্থাপনা (Human Resource Management বা HRM):} কর্মীর পরিচয়, উপস্থিতি, ছুটি, পদ ও প্রশিক্ষণের তথ্য পরিচালনা করে। এ কাজে ব্যবহৃত সফটওয়্যার ব্যবস্থাকে Human Resource Management System বা HRMS বলা হয়।
\item \textbf{হিসাব ও অর্থব্যবস্থা (Accounting and Finance):} আয়-ব্যয়, পাওনা, দেনা ও আর্থিক হিসাব রাখে। এসব তথ্য সংগ্রহ, সংরক্ষণ ও প্রক্রিয়াকরণে ব্যবহৃত ব্যবস্থাকে \textbf{Accounting Information System (AIS)} বলে। একই EmployeeID দিয়ে মানবসম্পদ বিভাগ কর্মীর পরিচয় এবং হিসাব বিভাগ বেতনের তথ্য ব্যবহার করতে পারে। প্রত্যেক বিভাগকে সব ব্যক্তিগত তথ্য দেখানোর প্রয়োজন নেই।
\item \textbf{মজুদ ব্যবস্থাপনা (Inventory Management):} পণ্যের মজুদ, ক্রয় ও বিক্রয়ের পরিমাণ অনুসরণ করে।
\item \textbf{সরবরাহ ব্যবস্থা পরিচালনা (Supply Chain Management বা SCM):} সরবরাহকারী, সংগ্রহ, পরিবহন ও বিতরণের তথ্য সমন্বয় করে। উৎপাদন (Manufacturing) অংশে কাঁচামাল ও উৎপাদন এবং বিতরণ (Distribution) অংশে গুদাম থেকে পণ্য পৌঁছানোর তথ্য থাকে। একটি পণ্য বিক্রি হলে মজুদ কমা এবং প্রয়োজন হলে নতুন ক্রয়ের সংকেত পাওয়ার জন্য এসব তথ্য সংযুক্ত থাকে।
\item \textbf{গ্রাহক সম্পর্ক ব্যবস্থাপনা (Customer Relationship Management বা CRM):} গ্রাহকের পরিচয়, ক্রয়-ইতিহাস, অভিযোগ ও সেবার তথ্য ব্যবহার করে সম্পর্ক ও সেবা পরিচালনা করে।
\item \textbf{এন্টারপ্রাইজ রিসোর্স প্ল্যানিং (Enterprise Resource Planning বা ERP):} প্রতিষ্ঠানের বিভিন্ন কার্যক্রমকে সমন্বিত ব্যবস্থায় আনে। ERP শুধু একটি টেবিল বা একটি DBMS-এর নাম নয়; এর অ্যাপ্লিকেশনগুলো ডেটাবেজের তথ্য ব্যবহার করে কাজ সমন্বয় করে।
\end{enumerate}

\subsection{একটি সমন্বিত উদাহরণ}
\par\smallskip\noindent\begin{minipage}[t]{.53\linewidth}\vspace{0pt}\raggedright
\small
\begin{tabularx}{\linewidth}{|p{22mm}|X|}\hline
\chSixHeader ব্যবস্থা & একই অর্ডারের তথ্য ব্যবহার\\\hline
Sales/CRM & গ্রাহক, চাহিদা ও অর্ডার\\\hline
Inventory & মজুদ যাচাই ও সরবরাহের পর পরিমাণ পরিবর্তন\\\hline
Accounts & মূল্য, বিল ও পরিশোধ\\\hline
Distribution & গন্তব্য ও পৌঁছানোর অবস্থা\\\hline
Management & বিক্রয়, ব্যয় ও কার্যক্রমের সারাংশ\\\hline
\end{tabularx}
\end{minipage}\hfill
\begin{minipage}[t]{.44\linewidth}\vspace{0pt}\raggedright
একটি প্রতিষ্ঠানে গ্রাহক ২০টি চেয়ার কিনলেন। বিক্রয় বিভাগ অর্ডার লিখল, মজুদ বিভাগ সরবরাহযোগ্য চেয়ার যাচাই করল, হিসাব বিভাগ বিল তৈরি করল এবং পরিবহন বিভাগ গন্তব্য পেল। একই OrderID ধরে কাজগুলো যুক্ত থাকলে বারবার নতুন করে গ্রাহকের পরিচয় লিখতে হয় না। অর্ডার বাতিল বা পরিমাণ বদল হলে সংশ্লিষ্ট বিভাগ সঠিক হালনাগাদ পেতে পারে।
\end{minipage}\par\smallskip


\subsection{স্থাপনার ধরন}
\begin{enumerate}
\item \textbf{Client--Server:} ব্যবহারকারীর অ্যাপ্লিকেশন কেন্দ্রীয় সার্ভারের ডেটা ব্যবহার করে।
\item \textbf{Distributed:} তথ্য বা তার অংশ বিভিন্ন স্থানে থাকতে পারে।
\item \textbf{Web-based:} ব্রাউজারের মাধ্যমে সেবা পাওয়া যায়।
\end{enumerate}
\textbf{মনে রাখবে:} এগুলো ব্যবহার ও স্থাপনার ধরন; রিলেশনাল মডেলের সরাসরি বিকল্প শ্রেণি নয়।

\subsubsection{সুবিধা}
\begin{enumerate}
\item \textbf{সমন্বয়:} বিভিন্ন বিভাগের তথ্য যুক্ত থাকলে প্রতিবেদন তৈরি, পুনরাবৃত্তি কমানো ও কাজের সামঞ্জস্য রাখা সহজ হয়।
\item \textbf{প্রয়োগের পরিসর:} ছোট প্রতিষ্ঠানেও এই নীতি প্রয়োগ করা যায়।
\end{enumerate}
\subsubsection{সীমাবদ্ধতা ও নির্বাচনের শর্ত}
\begin{enumerate}
\item \textbf{প্রস্তুতি ও রক্ষণাবেক্ষণ:} সঠিক নকশা, অনুমতি, প্রশিক্ষণ, নির্ভরযোগ্য নেটওয়ার্ক ও ব্যাকআপ প্রয়োজন।
\item \textbf{উপযুক্ত সফটওয়্যার:} তথ্যের পরিমাণ, ব্যবহারকারীর সংখ্যা, লেনদেনের চাপ এবং পরিচালনার সামর্থ্য দেখে সফটওয়্যার নির্বাচন করতে হয়।
\item \textbf{ভুল তথ্যের ঝুঁকি:} শুধু বড় ডেটাবেজ ব্যবহার করলে ভুল তথ্য বা দুর্বল কাজের পদ্ধতি ঠিক হয়ে যায় না।
\end{enumerate}

\subsection{কর্পোরেট কার্যক্রমের পূর্ণ চক্র}
একটি পোশাক প্রতিষ্ঠান গ্রাহকের অর্ডার পেয়ে কাঁচামাল সংগ্রহ, উৎপাদন, মান যাচাই, মজুদ, পরিবহন ও অর্থপ্রাপ্তির কাজ করে। ক্রয় বিভাগ SupplierID, উৎপাদন বিভাগ BatchID এবং বিক্রয় বিভাগ OrderID ব্যবহার করতে পারে। এসব পরিচিতির মাধ্যমে সংশ্লিষ্ট তথ্য যুক্ত থাকলে কোন অর্ডারের জন্য কোন কাঁচামাল ব্যবহৃত হয়েছে এবং কোন পণ্য পাঠানো হয়েছে জানা যায়।

মানবসম্পদ বিভাগ শ্রমিকের উপস্থিতি রাখে, হিসাব বিভাগ অনুমোদিত উপস্থিতি ও বেতনের নিয়ম ব্যবহার করে বেতন তৈরি করে। CRM গ্রাহকের অভিযোগ ধরে এবং সংশ্লিষ্ট অর্ডারের তথ্য দিয়ে সেবা দেয়। ব্যবস্থাপক পৃথক রেকর্ডের বদলে বিভাগভিত্তিক সারাংশ দেখতে পারেন। এ কারণে একই ডেটার বিভিন্ন উপস্থাপন ও ভিন্ন অনুমতি দরকার হয়।

ডেটাবেজের আকার বড় হলেই ব্যবস্থাপনা ভালো হবে না। একই পণ্যের জন্য দুই বিভাগে দুটি ভিন্ন কোড থাকলে সমন্বয় ব্যাহত হয়। পরিকল্পিত পরিচিতি, একই অর্থের ফিল্ড এবং যথাযথ হালনাগাদ প্রক্রিয়া প্রয়োজন। ছোট প্রতিষ্ঠানে সীমিত টেবিল দিয়েও এই নীতি প্রয়োগ করে পরে কাজের সঙ্গে ব্যবস্থা বড় করা যায়।


\textbf{সমন্বয় কোথায় কাজে লাগে?} বিক্রয় বিভাগ পণ্যের অর্ডার নিলেও মজুতের তথ্য না মিললে মজুতে নেই এমন পণ্য বিক্রির প্রতিশ্রুতি দিতে পারে। একই ProductID দিয়ে অর্ডার ও মজুত মিলিয়ে এবং পরিবর্তনের নিয়ম নিয়ন্ত্রণ করে এই ঝুঁকি কমে। বিভাগগুলোর তথ্যের সময়, পরিচয় ও অনুমোদনের নিয়ম মিলতে হবে; শুধু একই সার্ভারে আলাদা ফাইল রাখাকে পূর্ণ সমন্বয় বলা যায় না।

\subsection{বড় প্রতিষ্ঠানের Data সংরক্ষণের প্রয়োজন}
কর্পোরেট ডেটাবেজে অনেক বিভাগ, বিপুল রেকর্ড ও একাধিক ব্যবহারকারী থাকে। বিক্রয় চলার সময় হিসাব ও গুদাম বিভাগও তথ্য ব্যবহার করে; তাই তথ্যের সামঞ্জস্য, দ্রুত সাড়া, অনুমতি নিয়ন্ত্রণ, প্রয়োজনমতো সম্প্রসারণ এবং ব্যাকআপ গুরুত্বপূর্ণ। ক্রয় থেকে বিক্রয় পর্যন্ত একই পণ্যের পরিচিতি বজায় থাকলে বিভিন্ন বিভাগের হিসাব মিলিয়ে দেখা যায়।

দৈনন্দিন লেনদেনের পাশাপাশি দীর্ঘ সময়ের তথ্য বিশ্লেষণের প্রয়োজন হয়। বিভিন্ন উৎসের সংগঠিত তথ্য বিশ্লেষণের জন্য একত্রে রাখা হলে \textbf{Data Warehouse}-এর ধারণা পাওয়া যায়। বিভিন্ন ধরনের কাঁচা ডেটা বড় সংগ্রহে রাখার প্রসঙ্গে \textbf{Data Lake} ব্যবহৃত হয়। ভিন্ন ডেটা মডেলের NoSQL ব্যবস্থা ও বৃহৎ তথ্যের বিশ্লেষণ যথাক্রমে \ref{sec:ch06-nosql} ও \ref{sec:ch06-big-data} অংশে ব্যাখ্যা করা হয়েছে। যেমন নিয়মিত বিক্রয়ের টেবিল দিয়ে মাসের হিসাব এবং কয়েক বছরের সংগঠিত বিক্রয়তথ্য দিয়ে প্রবণতা বিশ্লেষণ ভিন্ন প্রয়োজন পূরণ করে।

\subsection{গুরুত্বপূর্ণ প্রশ্ন}
\textbf{অনুশীলন:} একটি দোকানে বিক্রয় লিখলে মজুদ ও পাওনার হিসাব একসঙ্গে বদলাতে হয় কেন? গ্রাহকের অভিযোগ কোন ব্যবস্থার তথ্য এবং কর্মীর ছুটি কোন ব্যবস্থার তথ্য---ব্যাখ্যা করো।


